\chapter{Lung Abscess}
\index{Lung Abscess}

\section*{Definition and Overview}
A lung abscess is a localized, cavitating area of pulmonary necrosis and suppuration (pus formation) resulting from a microbial infection. It is typically characterized by a single or multiple dominant cavities, often containing an air-fluid level on chest imaging. The primary mechanism of formation is the aspiration of oropharyngeal secretions containing anaerobic bacteria into the lower respiratory tract, particularly in patients with altered consciousness or impaired swallowing reflexes. Lung abscesses represent a severe form of necrotising pulmonary infection. While the introduction of modern antibiotics has significantly reduced the incidence and mortality of this condition, it remains a serious illness that requires prolonged antimicrobial therapy, careful diagnostic evaluation to rule out underlying malignancies, and occasionally, surgical or percutaneous intervention.

\section*{Epidemiology}
The epidemiology of lung abscesses has evolved substantially since the pre-antibiotic era, when it was a highly fatal condition. Today, it is relatively uncommon, primarily affecting individuals with predisposing risk factors for aspiration or compromised immune function. It occurs more frequently in males than in females and is more common in older adults due to a higher prevalence of dysphagia, neurological disorders, and periodontal disease. In many countries and other developed countries, the risk is highly correlated with alcohol dependence, drug overdose, epilepsy, stroke, and general anaesthesia. Poor oral health and severe periodontal disease are critical epidemiological links, as they significantly increase the bacterial load of anaerobic pathogens in the saliva.

\section*{Aetiology and Pathophysiology}
The development of a lung abscess is a multi-step process beginning with the entry of pathogenic organisms into the lung tissue, followed by tissue destruction and cavitation.
\begin{itemize}
    \item \textbf{Aspiration of oral secretions:} The most common cause. Aspiration occurs when protective airway reflexes are compromised, allowing saliva and oral bacteria to enter the bronchial tree, typically settling in the gravity-dependent segments of the lungs (the posterior segments of the upper lobes and apical segments of the lower lobes).
    \item \textbf{Anaerobic bacteria:} The most frequently isolated pathogens, including \textit{Peptostreptococcus} species, \textit{Prevotella} species, \textit{Fusobacterium nucleatum}, and \textit{Bacteroides} species. These bacteria are normal inhabitants of the gingival crevice.
    \item \textbf{Aerobic and opportunistic pathogens:} Less common but highly virulent causes include \textit{Staphylococcus aureus} (especially MRSA, which can cause necrotising pneumonia), \textit{Klebsiella pneumoniae} (common in patients with diabetes or alcohol overuse), \textit{Pseudomonas aeruginosa}, and \textit{Streptococcus pneumoniae}.
    \item \textbf{Bronchial obstruction:} An obstructed bronchus due to a lung tumour (bronchogenic carcinoma) or an inhaled foreign body prevents normal clearance of secretions, leading to post-obstructive pneumonia and abscess formation.
    \item \textbf{Haematogenous seeding:} Septic emboli traveling through the bloodstream from distant infections (such as tricuspid valve infective endocarditis or septic thrombophlebitis/Lemierre's syndrome) can lodge in the pulmonary vasculature and cause multi-focal abscesses.
\end{itemize}

\section*{Clinical Features}
The symptoms of a lung abscess typically develop subacutely over weeks, though acute presentations can occur with virulent aerobic pathogens.
\begin{itemize}
    \item \textbf{Productive cough:} Coughing up copious amounts of purulent sputum, which is characteristically foul-smelling and bad-tasting if anaerobic bacteria are involved.
    \item \textbf{Constitutional symptoms:} High fever, night sweats, rigors, malaise, anorexia, and progressive weight loss.
    \item \textbf{Chest pain:} Pleuritic chest pain that worsens with deep breathing or coughing, indicating inflammation of the adjacent pleura.
    \item \textbf{Dyspnoea:} Shortness of breath, particularly during exertion or when lying flat.
    \item \textbf{Haemoptysis:} Blood-streaked sputum or, less commonly, massive haemoptysis due to erosion of blood vessels within the abscess cavity.
    \item \textbf{Physical signs:} Dullness to percussion, bronchial breath sounds, and crackles over the affected lung area. Finger clubbing may develop in chronic, untreated cases.
\end{itemize}

\section*{Differential Diagnosis}
A cavitating lung lesion on imaging presents a diagnostic challenge, as several non-infectious and chronic infectious conditions can mimic a lung abscess.
\begin{itemize}
    \item \textbf{Cavitating lung malignancy:} Squamous cell carcinoma and other lung cancers can undergo central necrosis and cavitation, presenting with fever and cough, and must be ruled out in all patients, especially smokers.
    \item \textit{Mycobacterium tuberculosis} infection classically presents with cavitary lesions in the upper lobes, accompanied by weight loss, night sweats, and haemoptysis.
    \item \textbf{Necrotising pneumonia:} Characterised by multiple small cavities ($< 2$~cm) throughout the lung parenchyma rather than a single, large abscess cavity.
    \item \textbf{Cavitating pulmonary infarction:} A pulmonary embolism leading to lung tissue death and secondary aseptic or septic cavitation.
    \item \textbf{Fungal infections:} Cavitary lesions caused by \textit{Aspergillus} species (including aspergilloma), \textit{Coccidioides}, or \textit{Histoplasma}, typically in immunocompromised hosts.
    \item \textbf{Wegener's granulomatosis (GPA):} An autoimmune vasculitis that can cause necrotising granulomas and cavitating nodules in the lungs.
\end{itemize}

\section*{When to Seek Medical Care}
Patients presenting with a persistent, worsening cough, high fevers, night sweats, or the production of foul-smelling sputum should seek urgent medical evaluation.

\textbf{Contact your local emergency services for:}
\begin{itemize}
    \item Coughing up large volumes of bright red blood (massive haemoptysis), which can obstruct the airway
    \item Severe, acute shortness of breath or blue discolouration of the lips and skin (cyanosis)
    \item High fever accompanied by confusion, rapid breathing, and low blood pressure (signs of septic shock)
    \item Sudden, sharp, severe chest pain associated with collapse
    \item Inability to clear secretions, leading to choking or airway compromise
\end{itemize}

\section*{Diagnosis and Investigations}
Diagnosis is established through chest imaging, supplemented by microbiological and bronchoscopic evaluations to determine the causative pathogen and rule out obstruction.
\begin{itemize}
    \item \textbf{Chest radiograph:} The initial screening tool, typically demonstrating a thick-walled cavity with an air-fluid level, surrounded by pulmonary consolidation.
    \item \textbf{Computed tomography (CT) of the chest:} Provides detailed cross-sectional views, allowing differentiation between a lung abscess and an empyema (pus in the pleural space). It also helps identify underlying bronchial tumours or foreign bodies.
    \item \textbf{Sputum testing:} Sputum Gram stain and culture, along with acid-fast bacilli (AFB) smears and cultures to rule out tuberculosis. Sputum cytology is performed if malignancy is suspected.
    \item \textbf{Blood cultures:} Obtained in febrile patients to detect bacteraemia and confirm haematogenous origin.
    \item \textbf{Bronchoscopy:} Indicated if a post-obstructive cause (such as a foreign body or bronchogenic carcinoma) is suspected. It allows for direct inspection of the airways and the collection of uncontaminated specimens via bronchoalveolar lavage (BAL) or protected brush catheters.
    \item \textbf{Blood tests:} Full blood count demonstrating marked leukocytosis with a neutrophil predominance, and elevated inflammatory markers (CRP and ESR).
\end{itemize}

\section*{Management}
The cornerstone of lung abscess treatment is prolonged antibiotic therapy combined with measures to facilitate the drainage of secretions.
\begin{itemize}
    \item \textbf{Antimicrobial therapy:} Initiated intravenously and transitioned to oral therapy once the patient is clinically stable. First-line choices include beta-lactam/beta-lactamase inhibitor combinations (e.g., amoxicillin-clavulanate or piperacillin-tazobactam) or clindamycin to ensure coverage of both oral anaerobes and streptococci.
    \item \textbf{Duration of treatment:} Antibiotics are typically continued for 3 to 6 weeks, or until chest imaging demonstrates complete resolution of the cavity or leaves only a small, stable scar.
    \item \textbf{Postural drainage and chest physiotherapy:} Encouraged to help clear the thick, purulent secretions from the abscess cavity, provided the airway is protected.
    \item \textbf{Bronchoscopic intervention:} Performed to retrieve aspirated foreign bodies or to aspirate thick mucus plugs obstructing the bronchus.
    \item \textbf{Percutaneous or surgical drainage:} Indicated in less than 10\% of cases. It is reserved for patients who fail to respond to antibiotics, have giant cavities ($> 6$~cm) at risk of rupture, or develop a concurrent empyema.
    \item \textbf{Surgical resection (lobectomy):} A last resort for massive haemoptysis, suspected malignancy, or persistent disease despite adequate medical therapy.
\end{itemize}

\section*{Complications}
\begin{itemize}
    \item \textbf{Empyema:} Spread of the infection into the pleural space, leading to a collection of pus (often requiring chest tube drainage).
    \item \textbf{Bronchopleural fistula:} An abnormal communication between the bronchial tree and the pleural space, leading to pneumothorax and persistent infection.
    \item \textbf{Massive haemoptysis:} Life-threatening arterial bleeding within the cavity, leading to asphyxiation.
    \item \textbf{Haematogenous dissemination:} Spread of bacteria to the systemic circulation, potentially causing brain abscesses, meningitis, or metastatic osteomyelitis.
    \item \textbf{Respiratory failure:} Severe hypoxia and respiratory distress requiring mechanical ventilation.
\end{itemize}

\section*{Prognosis and Outcomes}
The prognosis for a lung abscess is generally favourable with modern antibiotic therapy, with cure rates exceeding 90\% for primary aspiration abscesses. The clinical response to treatment is typically slow; fevers may persist for 7 to 10 days after starting antibiotics, and complete radiological clearance of the cavity can take several weeks or months. A poor prognosis is associated with older age, severe pre-existing comorbidities, immunocompromised status, infection with multi-resistant aerobic pathogens (such as \textit{Pseudomonas aeruginosa} or MRSA), or when the abscess is secondary to an obstructing bronchogenic carcinoma.

\section*{Prevention and Self-Care}
\begin{itemize}
    \item \textbf{Maintain excellent oral hygiene:} Brush teeth twice daily, floss, and visit the dentist regularly to reduce the burden of anaerobic bacteria in the saliva.
    \item \textbf{Protect airways during altered consciousness:} Elevate the head of the bed to 30--45 degrees for patients with dysphagia, reflux, or those recovering from sedation.
    \item \textbf{Manage swallowing disorders:} Work with a speech pathologist for swallowing evaluation and therapy if you have a history of stroke or neurological disease.
    \item \textbf{Complete the full antibiotic course:} Do not stop taking prescribed antibiotics early, even if symptoms improve, to ensure complete eradication of the infection.
    \item \textbf{Quit smoking:} Stop smoking to restore normal mucociliary clearance in the airways and improve lung defense mechanisms.
    \item \textbf{Seek early treatment for pneumonia:} Consult a doctor promptly if you develop symptoms of a chest infection to prevent progression to necrosis.
\end{itemize}

\section*{Sources and Further Reading}
The following organisations provide reliable information and guidelines on lung health, respiratory infections, and lung abscesses.
\begin{itemize}
    \item Lung Foundation Australia. \textit{Support, education, and resources for respiratory conditions.} https://lungfoundation.com.au/
    \item Thoracic Society of Australia and New Zealand (TSANZ). \textit{Clinical guidelines and position statements.} https://www.thoracic.org.au/
    \item NHS. \textit{Lung abscess: overview and treatment.} https://www.nhs.uk/
    \item Mayo Clinic. \textit{Pneumonia and lung infections.} https://www.mayoclinic.org/
    \item MSD Manual. \textit{Lung abscess: professional evaluation and management.} https://www.msdmanuals.com/
\end{itemize}
